% ==============================================================================
%  PROTOKOLL JUGENDAUSSCHUSS LEONBERG
%  Kompilierbar mit: pdflatex / xelatex / lualatex / tectonic
% ==============================================================================

\documentclass[11pt,a4paper]{article}

% --- Grundlegende Pakete & Schriftart ---
\usepackage[utf8]{inputenc}
\usepackage[T1]{fontenc}
\usepackage[ngerman]{babel}
\usepackage[default]{sourcesanspro}
\renewcommand{\familydefault}{\sfdefault}

% --- Seitenlayout ---
\usepackage[
    a4paper,
    left=18mm,
    right=18mm,
    top=14mm,
    bottom=18mm,
    headheight=0mm,
    headsep=0mm,
    footskip=10mm
]{geometry}

% --- Farben ---
\usepackage[dvipsnames]{xcolor}
\definecolor{jaPrimary}{HTML}{4A227A}      % Edles Jugendausschuss-Violett
\definecolor{jaSecondary}{HTML}{6B3E9E}    % Sekundär-Violett
\definecolor{jaPillBg}{HTML}{9F7AEA}       % Sanfteres Violett für Sektionstitel
\definecolor{jaLightBg}{HTML}{F6F3FA}      % Sehr dezenter, edler Hintergrund für Sektionsleiste
\definecolor{jaBorder}{HTML}{CBD5E1}       % Feine Rahmenfarbe
\definecolor{darkText}{HTML}{1A202C}       % Textfarbe Anthrazit
\definecolor{subtleGray}{HTML}{718096}     % Dezentes Grau für Hinweise
\definecolor{subtleText}{HTML}{4A5568}     % Untertitel / Klammerzusätze

% --- Pakete für Grafik, Boxen & Tabellen ---
\usepackage{graphicx}
\usepackage{fancyhdr}
\usepackage{calc}
\usepackage{needspace}
\usepackage[skins,breakable]{tcolorbox}
\usepackage{lastpage}
\usepackage{hyperref}
\usepackage{tabularx}
\usepackage{booktabs}
\usepackage{array}
\usepackage{colortbl}
\usepackage{color}
\arrayrulecolor{jaBorder}

\hypersetup{
    colorlinks=false,
    pdfborder={0 0 0},
    pdfauthor={Jugendausschuss Leonberg},
    pdftitle={Protokoll Jugendausschuss}
}

% --- Standardabstände ---
\setlength{\parindent}{0pt}
\setlength{\parskip}{0pt}
\linespread{1.15}

% --- Box-Stile (Dynamische Höhe mit intelligentem Mindestplatz & dezenten Fortsetzungs-Ecken) ---
\newtcolorbox{contentbox}[1][]{
    enhanced,
    breakable,
    colback=white,
    colframe=jaBorder,
    boxrule=0.8pt,
    arc=2.5mm,
    top=3.5mm,
    bottom=3.5mm,
    left=4mm,
    right=4mm,
    pad at break=2mm,
    before=\needspace{4.0cm},
    overlay first={
        \draw[jaBorder, line width=0.8pt] (frame.south west) -- ++(4mm,0);
        \draw[jaBorder, line width=0.8pt] (frame.south east) -- ++(-4mm,0);
    },
    overlay last={
        \draw[jaBorder, line width=0.8pt] (frame.north west) -- ++(4mm,0);
        \draw[jaBorder, line width=0.8pt] (frame.north east) -- ++(-4mm,0);
    },
    overlay middle={
        \draw[jaBorder, line width=0.8pt] (frame.south west) -- ++(4mm,0);
        \draw[jaBorder, line width=0.8pt] (frame.south east) -- ++(-4mm,0);
        \draw[jaBorder, line width=0.8pt] (frame.north west) -- ++(4mm,0);
        \draw[jaBorder, line width=0.8pt] (frame.north east) -- ++(-4mm,0);
    },
    #1
}

% Titel der Boxen (Clean & Schwarz)
\newcommand{\sectionpill}[1]{{\bfseries\color{darkText} #1}}

% --- Kopf- und Fußzeile ---
\pagestyle{fancy}
\fancyhf{}
\renewcommand{\headrulewidth}{0pt}
\renewcommand{\footrulewidth}{0pt}
\fancyfoot[R]{\small\color{darkText} Seite \thepage/\pageref{LastPage}}

\begin{document}
\color{darkText}

% --- Eleganter Kopfbereich mit JA-Logo, Titelhierarchie & Zierlinie ---
\noindent
\begin{minipage}[c]{0.18\textwidth}
    \raggedright
    \IfFileExists{logo_ja.png}{%
        \includegraphics[height=2.0cm,keepaspectratio]{logo_ja.png}%
    }{%
        \fbox{\parbox[c][1.8cm][c]{2.6cm}{\centering\bfseries\color{jaPrimary} Logo JA}}%
    }
\end{minipage}%
\hfill
\begin{minipage}[c]{0.80\textwidth}
    \raggedright
    {\fontsize{19}{23}\selectfont\bfseries\color{jaPrimary} PROTOKOLL JUGENDAUSSCHUSS}\\[1.8mm]
    {\large\bfseries\color{jaSecondary} Jugendausschuss der Stadt Leonberg}\\[1mm]
    {\footnotesize\color{subtleGray} Protokollvorlage für Projektgruppen}
\end{minipage}

\vspace{3mm}
{\color{jaPrimary}\hrule height 1.2pt}
\vspace{4.5mm}

% --- Metadaten-Zeilen ---
{\large\textbf{Datum:}} {{DATE_PLACEHOLDER}}

\vspace{4mm}

{\large\textbf{Projektgruppe:}} {{PROJECT_GROUP_PLACEHOLDER}}

\vspace{5mm}

% --- Box 1: Anwesende Mitglieder (Dynamisch) ---
\begin{contentbox}
    \sectionpill{Anwesende Mitglieder:}
    \vspace{2mm}
    
    {{ATTENDEES_PLACEHOLDER}}
\end{contentbox}

\vspace{3.5mm plus 1mm minus 1.5mm}

% --- Box 2: Hauptthemen des Treffens (Dynamisch) ---
\begin{contentbox}
    \sectionpill{Hauptthemen des Treffens:}
    \vspace{2mm}
    
    {{TOPICS_PLACEHOLDER}}
\end{contentbox}

\vspace{3.5mm plus 1mm minus 1.5mm}

% --- Box 3: Zusammenfassung der wichtigsten besprochenen Punkte (Dynamisch) ---
\begin{contentbox}
    \sectionpill{Zusammenfassung der wichtigsten besprochenen Punkte:}
    \vspace{2mm}
    
    {{SUMMARY_PLACEHOLDER}}
\end{contentbox}

\vspace{3.5mm plus 1mm minus 1.5mm}

% --- Box 4: Nächste Schritte (Dynamisch) ---
\begin{contentbox}
    \sectionpill{Nächste Schritte - wer jeweils zuständig \& bis wann fertig?}
    \vspace{2mm}
    
    {{NEXT_STEPS_PLACEHOLDER}}
\end{contentbox}

\vspace{3.5mm plus 1mm minus 1.5mm}

% --- Box 5: Ist Unterstützung benötigt? (Dynamisch) ---
\begin{contentbox}
    \sectionpill{Ist Unterstützung benötigt? Wann und wie?}\\[1.5mm]
    {\small\color{subtleText}(von Sprecher*innen, Stadtjugendref, Verwaltung,\dots)}
    \vspace{2mm}
    
    {{SUPPORT_PLACEHOLDER}}
\end{contentbox}

\vspace{3.5mm plus 1mm minus 1.5mm}

% --- Box 6: Wer muss beim nächsten Treffen zwingend dabei sein? (Dynamisch) ---
\begin{contentbox}
    \sectionpill{Wer muss beim nächsten Treffen zwingend dabei sein?}\\[1.5mm]
    {\small\color{subtleText}(Sprecher*innen, Stadtjugendref, Kooperationspartner,\dots)}
    \vspace{2mm}
    
    {{NEXT_ATTENDEES_PLACEHOLDER}}
\end{contentbox}

\vspace{4mm plus 1mm minus 1.5mm}

% --- Nächstes Treffen & Unterschriften (Zusammenhalten) ---
\nopagebreak
\noindent
{\textbf{Ort \& Zeit des nächsten Treffens:}} {{NEXT_MEETING_PLACEHOLDER}}\\[1.5mm]
{\small\color{subtleText}(falls schon festgelegt, sonst bitte möglichst zeitnah nachschicken!)}

\vspace{4mm plus 1mm minus 1.5mm}
\nopagebreak

% --- Unterschriftsfelder ---
\noindent
\begin{minipage}[t]{0.48\textwidth}
    \begin{tcolorbox}[
        enhanced,
        colback=white,
        colframe=jaBorder,
        boxrule=0.8pt,
        arc=1.5mm,
        height=3.0cm,
        top=2.5mm, bottom=2mm, left=3mm, right=3mm
    ]
        {\footnotesize\color{subtleText} Name in Druckschrift: {{SPEAKER_NAME_PLACEHOLDER}}}
    \end{tcolorbox}
    \vspace{1.5mm}
    {\small\textbf{Unterschrift Gruppensprecher*in}}
\end{minipage}%
\hfill
\begin{minipage}[t]{0.48\textwidth}
    \begin{tcolorbox}[
        enhanced,
        colback=white,
        colframe=jaBorder,
        boxrule=0.8pt,
        arc=1.5mm,
        height=3.0cm,
        top=2.5mm, bottom=2mm, left=3mm, right=3mm
    ]
        {\footnotesize\color{subtleText} Name in Druckschrift: {{RECORDER_NAME_PLACEHOLDER}}}
    \end{tcolorbox}
    \vspace{1.5mm}
    {\small\textbf{Unterschrift Protokollant*in}}
\end{minipage}

\end{document}
